Data files are organised into {\em records} and {\em comments}, whose syntax is detailed next.

\section{Records}

Each record includes:
\begin{itemize}
\item a leading {\em keyword}, which identifies the information provided in the record
\item a number of {\em fields}. A field is either a real number (numeric field) or a string of characters (character field). There is at least one field
\item a {\em terminating semicolumn} (;), which indicates the end of the record.
\end{itemize}

\begin{quote}
\begin{tabbing}
\faHandPointRight ~ \= \rm The following record, with keyword LINE, specifies a transmission line:\\
\> \tt LINE A-B BUS\_A BUS\_B  3.0 30.0 150.0 1400.0 1 ;
\end{tabbing}
\end{quote} 

Inside the record, the keyword, the fields and the terminating semicolon are separated by (at least one) space(s). Anything after the semicolumn is ignored. The next record (or comment) starts with the next line.

\begin{quote}
\begin{tabbing}
\faHandPointRight ~ \= \rm Two incorrect versions of the above sample record, owing to missing spaces:\\
\> \tt LINE A-B BUS\_ABUS\_B  3.0 30.0 150.0 1400.0 1 ;\\
\> \tt LINE A-B BUS\_A BUS\_B  3.030.0 150.0 1400.0 1 ;
\end{tabbing}
\end{quote} 

Inside a data file, a record may span over multiple lines; the semicolumn indicates the end of the record. Spanning over several lines is highly recommended for records that include many fields. Note that, depending upon the text editor and its settings, a long record could appear truncated when displayed.  

\begin{quote}
\begin{tabbing}
\faHandPointRight ~ \= \rm An example of long record spanning over three lines: \\
\> \tt INJEC GFOL VSC1 A 1.0 1.0 0.0 0.0 0.005 0.15 1.02 \\
\> \tt  1200.0 0.005 0.15 0.573 6.0 0.0033 0.0333 10.0 0.002 -999.0 \\
\> \tt  10.0 0.1667 50.0 0.10 0.4 0.5 1.0 -1000. -2000. 1.001 1 ;\\
\end{tabbing}
\end{quote} 

The nature of each field, and the total number of fields are specified in this documentation. Some records have optional fields, which are always located at the end of the record.

\subsection{Constraints on numeric fields}

There is almost no constraint on numeric fields. They are written in free format. The notation is that of floating-point numbers in \textsc{Matlab}: with or without a dot, with or without an exponent, exponent denoted by {\tt E} or {\tt D}.

\begin{quote}
\faHandPointRight ~ \rm Different valid formats of the same number: ~ \mbox{\tt 30 30.~30.0 3E01 3.E01 3.0E01 3.E1 3.e1}
\end{quote} 

\subsection{Constraints on character fields}

Character fields are limited to 20 characters. Only the first 20 characters are read; the remaining of the string is just ignored, without warning. It is thus discouraged to have more than 20 characters in a field. 

Furthermore, some character fields are limited to eight characters; the remaining of the string is just ignored.

Uppercase letters are significant within a character field. Thus, two fields that differ by the uppercase/lowercase spelling are different.

If a character field includes a space or a slash (/), it must be enclosed with quotes ( ' or '' ). In between the two quotes, the leading spaces are significant while the trailing ones are ignored. Keywords do not include spaces; hence, the use of quotes is useless.

Because the semi-column (;) is a special character (used to indicate the end of a record), {\bf \large it must not be included in any  character field}, even within quotes.

\section{Comments}

There are three ways to insert comments in the data files:
\begin{enumerate}
\item a line in which the first non blank character is an exclamation mark (!). At most the first 130 characters after the ! are memorised and reproduced on output files
\item a line in which the first non blank character is the sharp character (\#): this line is simply ignored by the program
\item anything written after the semicolon that terminates a record is also ignored. This is convenient to store (short) comments next to a record, without interfering with the latter.
\end{enumerate}

\begin{quote}
\begin{tabbing}
\faHandPointRight ~ \= \rm Comments starting with \# can be used to identify the fields of a record: \\
\> \tt  \# ~~~~~~~~~~~ name bus FP  ~FQ ~P  ~Q  \= \tt SNOM  RS  ~~LLS LSR  ~RR ~~~LLR\\
\> \tt  INJEC INDMACH1  SM   ~~2 ~0.2 0.2 0.~0.~0.~0.031 0.1 3.2 0.018 0.180\\
\> \tt  \# \> \tt    ~H  ~~A ~~B ~~LF \\
\> \> \tt    0.7 0.5 0.0 0.6 ;
\end{tabbing}
\end{quote} 

Comments do not span over several lines. If several lines are needed, each of them must start with a ! or a \#.

Empty lines are just ignored.


\section{Sharing data between files}

Records may be distributed over an arbitrary number of data files, which will be read sequentially. The order in which the records are placed inside the data files does not matter. The order in which the files are read does not matter either.

For instance, a first data file may be devoted to network data, a second one to data for the initial power flow computation, a third one to the dynamic data of the components connected to the network and a last one to simulation control parameters. The second and third files, for instance, may be swapped in the list without any effect.



